% MCQ Questions for Chapter 06
\section{বহুনির্বাচনি প্রশ্ন (MCQ)}

\subsection*{\textbf{\underline{ডেটাবেস বেসিক}}}

\begin{multicols}{2}
\begin{enumerate}

\item Database কী? \textbf{[ঢাকা-২০২৩]}
\begin{enumerate}[label=\alph*)]
\item সংগঠিত ডেটার সংগ্রহ
\item এলোমেলো ডেটা
\item একটি সফটওয়্যার
\item একটি হার্ডওয়্যার
\end{enumerate}

\item DBMS এর পূর্ণরূপ কী? \textbf{[চট্টগ্রাম-২০২৩]}
\begin{enumerate}[label=\alph*)]
\item Database Management System
\item Data Base Management Software
\item Digital Base Management System
\item Data Block Management System
\end{enumerate}

\item DBMS এর সুবিধা কোনটি? \textbf{[রাজশাহী-২০২২]}
\begin{enumerate}[label=\alph*)]
\item Data redundancy কমায়
\item Data redundancy বাড়ায়
\item Data নষ্ট করে
\item Data লুকায়
\end{enumerate}

\item Data Redundancy কী? \textbf{[সিলেট-২০২৩]}
\begin{enumerate}[label=\alph*)]
\item ডেটার পুনরাবৃত্তি
\item ডেটা মুছে ফেলা
\item ডেটা সংরক্ষণ
\item ডেটা প্রদর্শন
\end{enumerate}

\item Data Integrity কী? \textbf{[বরিশাল-২০২২]}
\begin{enumerate}[label=\alph*)]
\item ডেটার সঠিকতা ও সামঞ্জস্য
\item ডেটার পরিমাণ
\item ডেটার গতি
\item ডেটার রং
\end{enumerate}

\item Data Security কী? \textbf{[দিনাজপুর-২০২৩]}
\begin{enumerate}[label=\alph*)]
\item ডেটার সুরক্ষা
\item ডেটা মুছে ফেলা
\item ডেটা কপি করা
\item ডেটা প্রদর্শন
\end{enumerate}

\item Popular DBMS কোনটি? \textbf{[কুমিল্লা-২০২২]}
\begin{enumerate}[label=\alph*)]
\item MySQL, Oracle, MS SQL Server
\item MS Word, Excel
\item Photoshop, Illustrator
\item Chrome, Firefox
\end{enumerate}

\item File System এর তুলনায় DBMS এর সুবিধা কী? \textbf{[যশোর-২০২৩]}
\begin{enumerate}[label=\alph*)]
\item Data sharing সহজ
\item Data sharing কঠিন
\item Data নষ্ট হয়
\item Data লুকানো যায়
\end{enumerate}

\end{enumerate}
\end{multicols}

\subsection*{\textbf{\underline{Database Model}}}

\begin{multicols}{2}
\begin{enumerate}[resume]

\item সবচেয়ে জনপ্রিয় Database Model কোনটি? \textbf{[ঢাকা-২০২২]}
\begin{enumerate}[label=\alph*)]
\item Relational Model
\item Hierarchical Model
\item Network Model
\item Object-Oriented Model
\end{enumerate}

\item Relational Database এ ডেটা কীভাবে সংরক্ষিত হয়? \textbf{[চট্টগ্রাম-২০২২]}
\begin{enumerate}[label=\alph*)]
\item Table আকারে
\item Tree আকারে
\item Graph আকারে
\item List আকারে
\end{enumerate}

\item Table এর অপর নাম কী? \textbf{[রাজশাহী-২০২৩]}
\begin{enumerate}[label=\alph*)]
\item Relation
\item Entity
\item Attribute
\item Tuple
\end{enumerate}

\item Table এর Row কে কী বলে? \textbf{[সিলেট-২০২২]}
\begin{enumerate}[label=\alph*)]
\item Tuple বা Record
\item Attribute
\item Field
\item Column
\end{enumerate}

\item Table এর Column কে কী বলে? \textbf{[বরিশাল-২০২৩]}
\begin{enumerate}[label=\alph*)]
\item Attribute বা Field
\item Tuple
\item Record
\item Row
\end{enumerate}

\item Primary Key কী? \textbf{[দিনাজপুর-২০২২]}
\begin{enumerate}[label=\alph*)]
\item Unique identifier
\item Duplicate data
\item Foreign key
\item Index
\end{enumerate}

\item Foreign Key কী? \textbf{[কুমিল্লা-২০২৩]}
\begin{enumerate}[label=\alph*)]
\item অন্য table এর primary key
\item নিজের primary key
\item Unique key
\item Composite key
\end{enumerate}

\item Composite Key কী? \textbf{[যশোর-২০২২]}
\begin{enumerate}[label=\alph*)]
\item একাধিক field এর সমন্বয়ে key
\item একটি field এর key
\item Foreign key
\item Primary key
\end{enumerate}

\end{enumerate}
\end{multicols}

\subsection*{\textbf{\underline{SQL বেসিক}}}

\begin{multicols}{2}
\begin{enumerate}[resume]

\item SQL এর পূর্ণরূপ কী? \textbf{[ঢাকা-২০২৩]}
\begin{enumerate}[label=\alph*)]
\item Structured Query Language
\item Standard Query Language
\item Simple Query Language
\item System Query Language
\end{enumerate}

\item Database তৈরির command কোনটি? \textbf{[চট্টগ্রাম-২০২৩]}
\begin{enumerate}[label=\alph*)]
\item CREATE DATABASE
\item MAKE DATABASE
\item NEW DATABASE
\item BUILD DATABASE
\end{enumerate}

\item Table তৈরির command কোনটি? \textbf{[রাজশাহী-২০২২]}
\begin{enumerate}[label=\alph*)]
\item CREATE TABLE
\item MAKE TABLE
\item NEW TABLE
\item BUILD TABLE
\end{enumerate}

\item Data insert করার command কোনটি? \textbf{[সিলেট-২০২৩]}
\begin{enumerate}[label=\alph*)]
\item INSERT INTO
\item ADD INTO
\item PUT INTO
\item WRITE INTO
\end{enumerate}

\item Data select করার command কোনটি? \textbf{[বরিশাল-২০২২]}
\begin{enumerate}[label=\alph*)]
\item SELECT
\item GET
\item FETCH
\item READ
\end{enumerate}

\item সব column select করার syntax কোনটি? \textbf{[দিনাজপুর-২০২৩]}
\begin{enumerate}[label=\alph*)]
\item SELECT * FROM table
\item SELECT ALL FROM table
\item GET * FROM table
\item FETCH ALL FROM table
\end{enumerate}

\item Data update করার command কোনটি? \textbf{[কুমিল্লা-২০২২]}
\begin{enumerate}[label=\alph*)]
\item UPDATE
\item MODIFY
\item CHANGE
\item ALTER
\end{enumerate}

\item Data delete করার command কোনটি? \textbf{[যশোর-২০২৩]}
\begin{enumerate}[label=\alph*)]
\item DELETE
\item REMOVE
\item DROP
\item ERASE
\end{enumerate}

\end{enumerate}
\end{multicols}

\subsection*{\textbf{\underline{SQL Clauses}}}

\begin{multicols}{2}
\begin{enumerate}[resume]

\item WHERE clause এর কাজ কী? \textbf{[ঢাকা-২০২২]}
\begin{enumerate}[label=\alph*)]
\item শর্ত নির্ধারণ
\item সাজানো
\item গ্রুপিং
\item যোগ করা
\end{enumerate}

\item ORDER BY clause এর কাজ কী? \textbf{[চট্টগ্রাম-২০২২]}
\begin{enumerate}[label=\alph*)]
\item সাজানো
\item শর্ত নির্ধারণ
\item গ্রুপিং
\item যোগ করা
\end{enumerate}

\item ASC কী বোঝায়? \textbf{[রাজশাহী-২০২৩]}
\begin{enumerate}[label=\alph*)]
\item Ascending (ছোট থেকে বড়)
\item Descending (বড় থেকে ছোট)
\item Alphabetical
\item Numerical
\end{enumerate}

\item DESC কী বোঝায়? \textbf{[সিলেট-২০২২]}
\begin{enumerate}[label=\alph*)]
\item Descending (বড় থেকে ছোট)
\item Ascending (ছোট থেকে বড়)
\item Description
\item Default
\end{enumerate}

\item GROUP BY clause এর কাজ কী? \textbf{[বরিশাল-২০২৩]}
\begin{enumerate}[label=\alph*)]
\item একই মানের data গ্রুপ করা
\item সাজানো
\item শর্ত নির্ধারণ
\item যোগ করা
\end{enumerate}

\item HAVING clause কখন ব্যবহৃত হয়? \textbf{[দিনাজপুর-২০২২]}
\begin{enumerate}[label=\alph*)]
\item GROUP BY এর সাথে শর্ত দিতে
\item শুধু শর্ত দিতে
\item সাজাতে
\item যোগ করতে
\end{enumerate}

\item DISTINCT keyword এর কাজ কী? \textbf{[কুমিল্লা-২০২৩]}
\begin{enumerate}[label=\alph*)]
\item Duplicate মুছে ফেলা
\item সাজানো
\item গ্রুপিং
\item যোগ করা
\end{enumerate}

\item LIMIT clause এর কাজ কী? \textbf{[যশোর-২০২২]}
\begin{enumerate}[label=\alph*)]
\item নির্দিষ্ট সংখ্যক row দেখানো
\item সব row দেখানো
\item row মুছে ফেলা
\item row যোগ করা
\end{enumerate}

\end{enumerate}
\end{multicols}

\subsection*{\textbf{\underline{SQL Functions}}}

\begin{multicols}{2}
\begin{enumerate}[resume]

\item COUNT() function এর কাজ কী? \textbf{[ঢাকা-২০২৩]}
\begin{enumerate}[label=\alph*)]
\item Row গণনা করা
\item যোগ করা
\item গড় বের করা
\item সর্বোচ্চ বের করা
\end{enumerate}

\item SUM() function এর কাজ কী? \textbf{[চট্টগ্রাম-২০২৩]}
\begin{enumerate}[label=\alph*)]
\item যোগফল বের করা
\item গণনা করা
\item গড় বের করা
\item সর্বোচ্চ বের করা
\end{enumerate}

\item AVG() function এর কাজ কী? \textbf{[রাজশাহী-২০২২]}
\begin{enumerate}[label=\alph*)]
\item গড় বের করা
\item যোগ করা
\item গণনা করা
\item সর্বোচ্চ বের করা
\end{enumerate}

\item MAX() function এর কাজ কী? \textbf{[সিলেট-২০২৩]}
\begin{enumerate}[label=\alph*)]
\item সর্বোচ্চ মান বের করা
\item সর্বনিম্ন মান বের করা
\item গড় বের করা
\item যোগ করা
\end{enumerate}

\item MIN() function এর কাজ কী? \textbf{[বরিশাল-২০২২]}
\begin{enumerate}[label=\alph*)]
\item সর্বনিম্ন মান বের করা
\item সর্বোচ্চ মান বের করা
\item গড় বের করা
\item যোগ করা
\end{enumerate}

\item UPPER() function এর কাজ কী? \textbf{[দিনাজপুর-২০২৩]}
\begin{enumerate}[label=\alph*)]
\item বড় হাতের অক্ষরে রূপান্তর
\item ছোট হাতের অক্ষরে রূপান্তর
\item সংখ্যায় রূপান্তর
\item তারিখে রূপান্তর
\end{enumerate}

\item LOWER() function এর কাজ কী? \textbf{[কুমিল্লা-২০২২]}
\begin{enumerate}[label=\alph*)]
\item ছোট হাতের অক্ষরে রূপান্তর
\item বড় হাতের অক্ষরে রূপান্তর
\item সংখ্যায় রূপান্তর
\item তারিখে রূপান্তর
\end{enumerate}

\item LENGTH() function এর কাজ কী? \textbf{[যশোর-২০২৩]}
\begin{enumerate}[label=\alph*)]
\item String এর দৈর্ঘ্য বের করা
\item String যোগ করা
\item String মুছে ফেলা
\item String সাজানো
\end{enumerate}

\end{enumerate}
\end{multicols}

\subsection*{\textbf{\underline{SQL Joins}}}

\begin{multicols}{2}
\begin{enumerate}[resume]

\item JOIN কী? \textbf{[ঢাকা-২০২২]}
\begin{enumerate}[label=\alph*)]
\item একাধিক table থেকে data নেওয়া
\item একটি table থেকে data নেওয়া
\item Data মুছে ফেলা
\item Data যোগ করা
\end{enumerate}

\item INNER JOIN কী করে? \textbf{[চট্টগ্রাম-২০২২]}
\begin{enumerate}[label=\alph*)]
\item উভয় table এ মিল আছে এমন data
\item প্রথম table এর সব data
\item দ্বিতীয় table এর সব data
\item সব table এর সব data
\end{enumerate}

\item LEFT JOIN কী করে? \textbf{[রাজশাহী-২০২৩]}
\begin{enumerate}[label=\alph*)]
\item বাম table এর সব data
\item ডান table এর সব data
\item উভয় table এ মিল আছে এমন data
\item কোনো data নয়
\end{enumerate}

\item RIGHT JOIN কী করে? \textbf{[সিলেট-২০২২]}
\begin{enumerate}[label=\alph*)]
\item ডান table এর সব data
\item বাম table এর সব data
\item উভয় table এ মিল আছে এমন data
\item কোনো data নয়
\end{enumerate}

\item FULL OUTER JOIN কী করে? \textbf{[বরিশাল-২০২৩]}
\begin{enumerate}[label=\alph*)]
\item উভয় table এর সব data
\item শুধু বাম table এর data
\item শুধু ডান table এর data
\item মিল আছে এমন data
\end{enumerate}

\item CROSS JOIN কী করে? \textbf{[দিনাজপুর-২০২২]}
\begin{enumerate}[label=\alph*)]
\item Cartesian Product তৈরি করে
\item শুধু মিল আছে এমন data
\item শুধু বাম table এর data
\item শুধু ডান table এর data
\end{enumerate}

\item SELF JOIN কী? \textbf{[কুমিল্লা-২০২৩]}
\begin{enumerate}[label=\alph*)]
\item একই table নিজের সাথে join
\item দুটি ভিন্ন table join
\item তিনটি table join
\item কোনো join নয়
\end{enumerate}

\item JOIN করার সময় কোন clause ব্যবহৃত হয়? \textbf{[যশোর-২০২২]}
\begin{enumerate}[label=\alph*)]
\item ON
\item WHERE
\item HAVING
\item GROUP BY
\end{enumerate}

\end{enumerate}
\end{multicols}

\subsection*{\textbf{\underline{Database Normalization}}}

\begin{multicols}{2}
\begin{enumerate}[resume]

\item Normalization কী? \textbf{[ঢাকা-২০২৩]}
\begin{enumerate}[label=\alph*)]
\item Database সংগঠিত করার প্রক্রিয়া
\item Database তৈরির প্রক্রিয়া
\item Database মুছে ফেলার প্রক্রিয়া
\item Database backup করার প্রক্রিয়া
\end{enumerate}

\item Normalization এর উদ্দেশ্য কী? \textbf{[চট্টগ্রাম-২০২৩]}
\begin{enumerate}[label=\alph*)]
\item Redundancy কমানো
\item Redundancy বাড়ানো
\item Data মুছে ফেলা
\item Data লুকানো
\end{enumerate}

\item 1NF এর শর্ত কী? \textbf{[রাজশাহী-২০২২]}
\begin{enumerate}[label=\alph*)]
\item প্রতিটি cell এ একটি মাত্র মান
\item প্রতিটি cell এ একাধিক মান
\item Primary key থাকতে হবে
\item Foreign key থাকতে হবে
\end{enumerate}

\item 2NF এর শর্ত কী? \textbf{[সিলেট-২০২৩]}
\begin{enumerate}[label=\alph*)]
\item 1NF এবং Partial Dependency নেই
\item শুধু 1NF
\item শুধু Primary key
\item শুধু Foreign key
\end{enumerate}

\item 3NF এর শর্ত কী? \textbf{[বরিশাল-২০২২]}
\begin{enumerate}[label=\alph*)]
\item 2NF এবং Transitive Dependency নেই
\item শুধু 2NF
\item শুধু 1NF
\item কোনো শর্ত নেই
\end{enumerate}

\item Functional Dependency কী? \textbf{[দিনাজপুর-২০২৩]}
\begin{enumerate}[label=\alph*)]
\item একটি attribute অন্যটির উপর নির্ভরশীল
\item সব attribute স্বাধীন
\item কোনো নির্ভরশীলতা নেই
\item সব attribute একই
\end{enumerate}

\item Partial Dependency কী? \textbf{[কুমিল্লা-২০২২]}
\begin{enumerate}[label=\alph*)]
\item Composite key এর অংশের উপর নির্ভরশীল
\item পুরো key এর উপর নির্ভরশীল
\item কোনো নির্ভরশীলতা নেই
\item Foreign key এর উপর নির্ভরশীল
\end{enumerate}

\item Transitive Dependency কী? \textbf{[যশোর-২০২৩]}
\begin{enumerate}[label=\alph*)]
\item A→B এবং B→C হলে A→C
\item সরাসরি নির্ভরশীলতা
\item কোনো নির্ভরশীলতা নেই
\item Foreign key নির্ভরশীলতা
\end{enumerate}

\end{enumerate}
\end{multicols}

\subsection*{\textbf{\underline{উত্তরমালা}}}

\begin{multicols}{4}
\begin{enumerate}
\item a
\item a
\item a
\item a
\item a
\item a
\item a
\item a
\item a
\item a
\item a
\item a
\item a
\item a
\item a
\item a
\item a
\item a
\item a
\item a
\item a
\item a
\item a
\item a
\item a
\item a
\item a
\item a
\item a
\item a
\item a
\item a
\item a
\item a
\item a
\item a
\item a
\item a
\item a
\item a
\item a
\item a
\item a
\item a
\item a
\item a
\item a
\item a
\item a
\item a
\item a
\item a
\item a
\item a
\item a
\item a
\end{enumerate}
\end{multicols}

\subsection*{\textbf{\underline{SQL Advanced}}}

\begin{multicols}{2}
\begin{enumerate}[resume]

\item ALTER TABLE command কী করে? \textbf{[ঢাকা-২০২৩]}
\begin{enumerate}[label=\alph*)]
\item Table এর structure পরিবর্তন
\item Table তৈরি
\item Table মুছে ফেলা
\item Data insert
\end{enumerate}

\item DROP TABLE command কী করে? \textbf{[চট্টগ্রাম-২০২২]}
\begin{enumerate}[label=\alph*)]
\item Table মুছে ফেলা
\item Table তৈরি
\item Table পরিবর্তন
\item Data insert
\end{enumerate}

\item TRUNCATE command কী করে? \textbf{[সিলেট-২০২৩]}
\begin{enumerate}[label=\alph*)]
\item Table এর সব data মুছে ফেলা
\item Table মুছে ফেলা
\item Table তৈরি
\item Data insert
\end{enumerate}

\item DELETE এবং TRUNCATE এর পার্থক্য কী? \textbf{[রাজশাহী-২০২২]}
\begin{enumerate}[label=\alph*)]
\item DELETE শর্ত দিয়ে মুছে, TRUNCATE সব মুছে
\item কোনো পার্থক্য নেই
\item DELETE দ্রুত
\item TRUNCATE শর্ত দিয়ে মুছে
\end{enumerate}

\item INDEX কী? \textbf{[বরিশাল-২০২৩]}
\begin{enumerate}[label=\alph*)]
\item দ্রুত search এর জন্য structure
\item একটি table
\item একটি database
\item একটি column
\end{enumerate}

\item VIEW কী? \textbf{[দিনাজপুর-২০২২]}
\begin{enumerate}[label=\alph*)]
\item Virtual table
\item Real table
\item Database
\item Column
\end{enumerate}

\item TRIGGER কী? \textbf{[কুমিল্লা-২০২৩]}
\begin{enumerate}[label=\alph*)]
\item স্বয়ংক্রিয় চলা stored procedure
\item একটি table
\item একটি database
\item একটি column
\end{enumerate}

\item STORED PROCEDURE কী? \textbf{[যশোর-২০২২]}
\begin{enumerate}[label=\alph*)]
\item সংরক্ষিত SQL code
\item একটি table
\item একটি database
\item একটি column
\end{enumerate}

\item TRANSACTION কী? \textbf{[ঢাকা-২০২৩]}
\begin{enumerate}[label=\alph*)]
\item একসাথে চলা operations এর সেট
\item একটি table
\item একটি database
\item একটি column
\end{enumerate}

\item COMMIT command কী করে? \textbf{[চট্টগ্রাম-২০২২]}
\begin{enumerate}[label=\alph*)]
\item Transaction সংরক্ষণ
\item Transaction বাতিল
\item Transaction শুরু
\item Transaction শেষ
\end{enumerate}

\item ROLLBACK command কী করে? \textbf{[সিলেট-২০২৩]}
\begin{enumerate}[label=\alph*)]
\item Transaction বাতিল
\item Transaction সংরক্ষণ
\item Transaction শুরু
\item Transaction শেষ
\end{enumerate}

\item SAVEPOINT কী? \textbf{[রাজশাহী-২০২২]}
\begin{enumerate}[label=\alph*)]
\item Transaction এর মধ্যে checkpoint
\item Transaction শুরু
\item Transaction শেষ
\item Transaction বাতিল
\end{enumerate}

\item ACID property এর A কী? \textbf{[বরিশাল-২০২৩]}
\begin{enumerate}[label=\alph*)]
\item Atomicity
\item Accuracy
\item Availability
\item Authentication
\end{enumerate}

\item ACID property এর C কী? \textbf{[দিনাজপুর-২০২২]}
\begin{enumerate}[label=\alph*)]
\item Consistency
\item Completeness
\item Correctness
\item Compatibility
\end{enumerate}

\item ACID property এর I কী? \textbf{[কুমিল্লা-২০২৩]}
\begin{enumerate}[label=\alph*)]
\item Isolation
\item Integration
\item Implementation
\item Identification
\end{enumerate}

\item ACID property এর D কী? \textbf{[যশোর-২০২২]}
\begin{enumerate}[label=\alph*)]
\item Durability
\item Dependency
\item Distribution
\item Documentation
\end{enumerate}

\end{enumerate}
\end{multicols}

\subsection*{\textbf{\underline{Database Design}}}

\begin{multicols}{2}
\begin{enumerate}[resume]

\item ER Diagram এর পূর্ণরূপ কী? \textbf{[ঢাকা-২০২৩]}
\begin{enumerate}[label=\alph*)]
\item Entity Relationship Diagram
\item Entity Record Diagram
\item Electronic Relationship Diagram
\item Electronic Record Diagram
\end{enumerate}

\item Entity কী? \textbf{[চট্টগ্রাম-২০২২]}
\begin{enumerate}[label=\alph*)]
\item বাস্তব জগতের object
\item একটি table
\item একটি column
\item একটি row
\end{enumerate}

\item Attribute কী? \textbf{[সিলেট-২০২৩]}
\begin{enumerate}[label=\alph*)]
\item Entity এর বৈশিষ্ট্য
\item একটি table
\item একটি database
\item একটি row
\end{enumerate}

\item Relationship কী? \textbf{[রাজশাহী-২০২২]}
\begin{enumerate}[label=\alph*)]
\item Entity গুলোর মধ্যে সম্পর্ক
\item একটি table
\item একটি column
\item একটি row
\end{enumerate}

\item One-to-One relationship কী? \textbf{[বরিশাল-২০২৩]}
\begin{enumerate}[label=\alph*)]
\item একটি entity অন্য একটির সাথে সম্পর্কিত
\item একটি entity অনেকগুলোর সাথে সম্পর্কিত
\item অনেকগুলো entity অনেকগুলোর সাথে সম্পর্কিত
\item কোনো সম্পর্ক নেই
\end{enumerate}

\item One-to-Many relationship কী? \textbf{[দিনাজপুর-২০২২]}
\begin{enumerate}[label=\alph*)]
\item একটি entity অনেকগুলোর সাথে সম্পর্কিত
\item একটি entity একটির সাথে সম্পর্কিত
\item অনেকগুলো entity অনেকগুলোর সাথে সম্পর্কিত
\item কোনো সম্পর্ক নেই
\end{enumerate}

\item Many-to-Many relationship কী? \textbf{[কুমিল্লা-২০২৩]}
\begin{enumerate}[label=\alph*)]
\item অনেকগুলো entity অনেকগুলোর সাথে সম্পর্কিত
\item একটি entity একটির সাথে সম্পর্কিত
\item একটি entity অনেকগুলোর সাথে সম্পর্কিত
\item কোনো সম্পর্ক নেই
\end{enumerate}

\item Cardinality কী? \textbf{[যশোর-২০২২]}
\begin{enumerate}[label=\alph*)]
\item Relationship এর সংখ্যা
\item Entity এর সংখ্যা
\item Attribute এর সংখ্যা
\item Table এর সংখ্যা
\end{enumerate}

\item Weak Entity কী? \textbf{[ঢাকা-২০২৩]}
\begin{enumerate}[label=\alph*)]
\item যার নিজস্ব primary key নেই
\item যার primary key আছে
\item যার foreign key নেই
\item যার কোনো key নেই
\end{enumerate}

\item Strong Entity কী? \textbf{[চট্টগ্রাম-২০২২]}
\begin{enumerate}[label=\alph*)]
\item যার নিজস্ব primary key আছে
\item যার primary key নেই
\item যার foreign key নেই
\item যার কোনো key নেই
\end{enumerate}

\item Composite Attribute কী? \textbf{[সিলেট-২০২৩]}
\begin{enumerate}[label=\alph*)]
\item একাধিক অংশে বিভক্ত attribute
\item একটি অংশের attribute
\item কোনো অংশ নেই
\item সব attribute একই
\end{enumerate}

\item Derived Attribute কী? \textbf{[রাজশাহী-২০২২]}
\begin{enumerate}[label=\alph*)]
\item অন্য attribute থেকে হিসাব করা
\item সরাসরি সংরক্ষিত
\item কোনো মান নেই
\item সব মান একই
\end{enumerate}

\item Multi-valued Attribute কী? \textbf{[বরিশাল-২০২৩]}
\begin{enumerate}[label=\alph*)]
\item একাধিক মান থাকতে পারে
\item একটি মাত্র মান
\item কোনো মান নেই
\item সব মান একই
\end{enumerate}

\item Key Attribute কী? \textbf{[দিনাজপুর-২০২২]}
\begin{enumerate}[label=\alph*)]
\item যা entity কে uniquely identify করে
\item যেকোনো attribute
\item কোনো বিশেষত্ব নেই
\item সব attribute একই
\end{enumerate}

\end{enumerate}
\end{multicols}

\subsection*{\textbf{\underline{Database Security ও Backup}}}

\begin{multicols}{2}
\begin{enumerate}[resume]

\item Database Security কেন প্রয়োজন? \textbf{[কুমিল্লা-২০২৩]}
\begin{enumerate}[label=\alph*)]
\item অননুমোদিত access রোধ করতে
\item দ্রুত access করতে
\item Data বাড়াতে
\item Data কমাতে
\end{enumerate}

\item Authentication কী? \textbf{[যশোর-২০২২]}
\begin{enumerate}[label=\alph*)]
\item ব্যবহারকারী যাচাই
\item Data যাচাই
\item Table যাচাই
\item Database যাচাই
\end{enumerate}

\item Authorization কী? \textbf{[ঢাকা-২০২৩]}
\begin{enumerate}[label=\alph*)]
\item অনুমতি প্রদান
\item ব্যবহারকারী যাচাই
\item Data যাচাই
\item Table যাচাই
\end{enumerate}

\item GRANT command কী করে? \textbf{[চট্টগ্রাম-২০২২]}
\begin{enumerate}[label=\alph*)]
\item অনুমতি দেয়
\item অনুমতি কেড়ে নেয়
\item Data insert করে
\item Data delete করে
\end{enumerate}

\item REVOKE command কী করে? \textbf{[সিলেট-২০২৩]}
\begin{enumerate}[label=\alph*)]
\item অনুমতি কেড়ে নেয়
\item অনুমতি দেয়
\item Data insert করে
\item Data delete করে
\end{enumerate}

\item Backup কী? \textbf{[রাজশাহী-২০২২]}
\begin{enumerate}[label=\alph*)]
\item Data এর copy সংরক্ষণ
\item Data মুছে ফেলা
\item Data পরিবর্তন
\item Data দেখা
\end{enumerate}

\item Full Backup কী? \textbf{[বরিশাল-২০২৩]}
\begin{enumerate}[label=\alph*)]
\item সম্পূর্ণ database এর backup
\item শুধু পরিবর্তিত data এর backup
\item কোনো backup নয়
\item অর্ধেক backup
\end{enumerate}

\item Incremental Backup কী? \textbf{[দিনাজপুর-২০২২]}
\begin{enumerate}[label=\alph*)]
\item শুধু পরিবর্তিত data এর backup
\item সম্পূর্ণ database এর backup
\item কোনো backup নয়
\item অর্ধেক backup
\end{enumerate}

\item Recovery কী? \textbf{[কুমিল্লা-২০২৩]}
\begin{enumerate}[label=\alph*)]
\item Backup থেকে data ফিরিয়ে আনা
\item Data মুছে ফেলা
\item Data পরিবর্তন
\item Data দেখা
\end{enumerate}

\item Encryption কী? \textbf{[যশোর-২০২২]}
\begin{enumerate}[label=\alph*)]
\item Data কে গোপন কোডে রূপান্তর
\item Data দেখা
\item Data মুছে ফেলা
\item Data পরিবর্তন
\end{enumerate}

\item Decryption কী? \textbf{[ঢাকা-২০২৩]}
\begin{enumerate}[label=\alph*)]
\item গোপন কোড থেকে data ফিরিয়ে আনা
\item Data গোপন করা
\item Data মুছে ফেলা
\item Data পরিবর্তন
\end{enumerate}

\item SQL Injection কী? \textbf{[চট্টগ্রাম-২০২২]}
\begin{enumerate}[label=\alph*)]
\item নিরাপত্তা আক্রমণ
\item Data insert
\item Data update
\item Data select
\end{enumerate}

\item Data Integrity Constraint কী? \textbf{[সিলেট-২০২৩]}
\begin{enumerate}[label=\alph*)]
\item Data এর সঠিকতা নিশ্চিত করার নিয়ম
\item Data মুছে ফেলার নিয়ম
\item Data দেখার নিয়ম
\item Data পরিবর্তনের নিয়ম
\end{enumerate}

\item NOT NULL constraint কী করে? \textbf{[রাজশাহী-২০২২]}
\begin{enumerate}[label=\alph*)]
\item NULL মান রোধ করে
\item NULL মান অনুমোদন করে
\item সব মান NULL করে
\item কোনো কাজ করে না
\end{enumerate}

\item UNIQUE constraint কী করে? \textbf{[বরিশাল-২০২৩]}
\begin{enumerate}[label=\alph*)]
\item Duplicate মান রোধ করে
\item Duplicate মান অনুমোদন করে
\item সব মান একই করে
\item কোনো কাজ করে না
\end{enumerate}

\item CHECK constraint কী করে? \textbf{[দিনাজপুর-২০২২]}
\begin{enumerate}[label=\alph*)]
\item শর্ত অনুযায়ী মান যাচাই
\item সব মান গ্রহণ
\item কোনো মান গ্রহণ না
\item কোনো কাজ করে না
\end{enumerate}

\item DEFAULT constraint কী করে? \textbf{[কুমিল্লা-২০২৩]}
\begin{enumerate}[label=\alph*)]
\item ডিফল্ট মান নির্ধারণ
\item সব মান মুছে ফেলা
\item সব মান পরিবর্তন
\item কোনো কাজ করে না
\end{enumerate}

\item AUTO\_INCREMENT কী করে? \textbf{[যশোর-২০২২]}
\begin{enumerate}[label=\alph*)]
\item স্বয়ংক্রিয়ভাবে মান বাড়ায়
\item স্বয়ংক্রিয়ভাবে মান কমায়
\item মান স্থির রাখে
\item কোনো কাজ করে না
\end{enumerate}

\end{enumerate}
\end{multicols}

\subsection*{\textbf{\underline{উত্তরমালা (৫৭-১০৬)}}}

\begin{multicols}{4}
\begin{enumerate}[start=57]
\item a
\item a
\item a
\item a
\item a
\item a
\item a
\item a
\item a
\item a
\item a
\item a
\item a
\item a
\item a
\item a
\item a
\item a
\item a
\item a
\item a
\item a
\item a
\item a
\item a
\item a
\item a
\item a
\item a
\item a
\item a
\item a
\item a
\item a
\item a
\item a
\item a
\item a
\item a
\item a
\item a
\item a
\item a
\item a
\item a
\item a
\item a
\item a
\item a
\item a
\end{enumerate}
\end{multicols}
